% Theorem C: Decision-Theoretic Verifier and Chow Abstention with Explicit Costs
\section{Theorem C: Decision-Theoretic Verifier and Chow's Rule}
\label{sec:thmC}

The manuscript (Claim C072; Sections XVI, XXI-C/D) asserted that the ARES decision rule forms a "classical Neyman--Pearson test" that is "decision-theoretically optimal". We provide the precise decision-theoretic derivation, identify the exact necessary and sufficient conditions for optimality, and formulate the three-way decision problem via Chow's rule with explicit operational costs.

\subsection{Formal Decision Model}
Let $Y \in \{0, 1\}$ denote the binary veracity state of an evaluated claim ($Y=1$ for True, $Y=0$ for False). An autonomous system selects an action from a ternary action space:
\begin{equation}
\mathcal{A} = \{ a_{\text{T}} \text{ (Assert True)}, \; a_{\text{F}} \text{ (Assert False)}, \; a_{\text{abstain}} \text{ (Abstain / Escalate to HITL)} \}.
\end{equation}
Let the cost matrix $C(a, y)$ be defined by:
\begin{align}
C(a_{\text{T}}, 1) &= 0, & C(a_{\text{T}}, 0) &= c_{\text{FP}} > 0 \quad \text{(Cost of verifying false content)}, \nonumber \\
C(a_{\text{F}}, 0) &= 0, & C(a_{\text{F}}, 1) &= c_{\text{FN}} > 0 \quad \text{(Cost of flagging authentic content)}, \nonumber \\
C(a_{\text{abstain}}, 1) &= c_a > 0, & C(a_{\text{abstain}}, 0) &= c_a > 0 \quad \text{(Cost of human review / analyst queue delay)}.
\end{align}

\begin{theorem}[Bayes-Optimal Decision Rule under Cost Matrix]
\label{thm:chow_rule}
Let $p(c) \triangleq \mathbb{P}(Y = 1 \mid c) \in [0, 1]$ be the well-calibrated posterior probability that claim $c$ is true, and let $\ell(c) \triangleq \ln \frac{p(c)}{1 - p(c)}$ denote the posterior log-odds.
\begin{enumerate}
    \item \textbf{Optimal Decision Rule:} The decision that minimizes the conditional expected Bayes risk $\mathcal{R}(a \mid c) = \mathbb{E}[C(a, Y) \mid c]$ is:
    \begin{equation}
    \label{eq:chow_prob_rule}
    a^*(c) = \begin{cases}
    a_{\text{T}} & \text{if } p(c) \ge 1 - \frac{c_a}{c_{\text{FP}}}, \\
    a_{\text{F}} & \text{if } p(c) \le \frac{c_a}{c_{\text{FN}}}, \\
    a_{\text{abstain}} & \text{if } \frac{c_a}{c_{\text{FN}}} < p(c) < 1 - \frac{c_a}{c_{\text{FP}}}.
    \end{cases}
    \end{equation}

    \item \textbf{Thresholds in Log-Odds Space:} Expressed on the log-odds scale $\ell(c)$, the optimal decision boundaries are:
    \begin{equation}
    \label{eq:chow_logodds_rule}
    \tau_1 \triangleq \ln \left( \frac{c_{\text{FP}} - c_a}{c_a} \right), \quad \tau_0 \triangleq \ln \left( \frac{c_a}{c_{\text{FN}} - c_a} \right),
    \end{equation}
    yielding the decision rule: assert True if $\ell(c) \ge \tau_1$, assert False if $\ell(c) \le \tau_0$, and escalate to HITL if $\tau_0 < \ell(c) < \tau_1$.

    \item \textbf{Non-Triviality Condition:} The abstention region $(\tau_0, \tau_1)$ is non-empty if and only if the human review cost is strictly strictly less than the harmonic mean cost of error:
    \begin{equation}
    \label{eq:nontrivial_abstain}
    c_a < \frac{c_{\text{FP}} \cdot c_{\text{FN}}}{c_{\text{FP}} + c_{\text{FN}}}.
    \end{equation}
    If $c_a \ge \frac{c_{\text{FP}} c_{\text{FN}}}{c_{\text{FP}} + c_{\text{FN}}}$, human escalation is cost-ineffective and the optimal policy collapses to a binary threshold.

    \item \textbf{Scope of Neyman--Pearson Optimality:}
    \begin{itemize}
        \item When evidence sources $e_1, \dots, e_K$ are \emph{conditionally independent} given $Y$, the posterior log-odds decomposes linearly as $\ell(c) = \ell_{\text{prior}} + \sum_{k=1}^K \operatorname{LLR}_k(e_k)$. In this case, for fixed type-I error rate $\alpha$, thresholding $\ell(c)$ yields the Neyman--Pearson optimal test.
        \item In general multi-modal settings where sources exhibit statistical dependence, a linear combination $\sum_{k=1}^K w_k \operatorname{LLR}_k$ optimized via logistic regression minimizes cross-entropy loss within the linear class, but is \emph{not} uniformly most powerful or Neyman--Pearson optimal across arbitrary non-linear joint likelihoods.
    \end{itemize}
\end{enumerate}
\end{theorem}

\begin{proof}
The expected costs of the three actions given posterior probability $p = \mathbb{P}(Y = 1 \mid c)$ are:
\begin{align}
\mathcal{R}(a_{\text{T}} \mid c) &= c_{\text{FP}} \mathbb{P}(Y = 0 \mid c) = c_{\text{FP}} (1 - p), \\
\mathcal{R}(a_{\text{F}} \mid c) &= c_{\text{FN}} \mathbb{P}(Y = 1 \mid c) = c_{\text{FN}} p, \\
\mathcal{R}(a_{\text{abstain}} \mid c) &= c_a.
\end{align}
Action $a_{\text{T}}$ is preferred over $a_{\text{abstain}}$ if and only if:
\begin{equation}
c_{\text{FP}}(1 - p) \le c_a \iff 1 - p \le \frac{c_a}{c_{\text{FP}}} \iff p \ge 1 - \frac{c_a}{c_{\text{FP}}}.
\end{equation}
Similarly, action $a_{\text{F}}$ is preferred over $a_{\text{abstain}}$ if and only if:
\begin{equation}
c_{\text{FN}} p \le c_a \iff p \le \frac{c_a}{c_{\text{FN}}}.
\end{equation}
Converting $p = \frac{e^\ell}{1 + e^\ell}$ to the log-odds threshold for $a_{\text{T}}$:
\begin{equation}
\frac{e^{\tau_1}}{1 + e^{\tau_1}} = 1 - \frac{c_a}{c_{\text{FP}}} = \frac{c_{\text{FP}} - c_a}{c_{\text{FP}}} \implies e^{\tau_1} = \frac{c_{\text{FP}} - c_a}{c_a} \implies \tau_1 = \ln\left( \frac{c_{\text{FP}} - c_a}{c_a} \right).
\end{equation}
For $a_{\text{F}}$:
\begin{equation}
\frac{e^{\tau_0}}{1 + e^{\tau_0}} = \frac{c_a}{c_{\text{FN}}} \implies e^{\tau_0} = \frac{c_a}{c_{\text{FN}} - c_a} \implies \tau_0 = \ln\left( \frac{c_a}{c_{\text{FN}} - c_a} \right).
\end{equation}
The abstention interval is non-empty if and only if $\frac{c_a}{c_{\text{FN}}} < 1 - \frac{c_a}{c_{\text{FP}}}$:
\begin{equation}
c_a\left( \frac{1}{c_{\text{FN}}} + \frac{1}{c_{\text{FP}}} \right) < 1 \iff c_a < \frac{1}{\frac{1}{c_{\text{FP}}} + \frac{1}{c_{\text{FN}}}} = \frac{c_{\text{FP}} c_{\text{FN}}}{c_{\text{FP}} + c_{\text{FN}}}.
\end{equation}
\end{proof}
